\section{Thiết kế Monitoring và Incident}

Hệ thống giám sát (Monitoring) và phát hiện sự cố (Incident Detection) được thiết kế nhằm theo dõi toàn diện trạng thái hoạt động của từng Edge Node từ mức phần cứng, tiến trình thực thi đến chất lượng suy luận của mô hình AI.

\subsection{Không gian đặc trưng Telemetry (8 Telemetry Features)}

Thay vì chỉ giám sát tài nguyên máy chủ thông thường, nền tảng thu thập bộ 8 đặc trưng telemetry cốt lõi thuộc ba nhóm thực thể tại Edge Node: phần cứng (\textit{Hardware}), môi trường thực thi (\textit{Runtime}) và chất lượng đầu ra mô hình (\textit{AI Output}), như mô tả trong Bảng~\ref{tab:telemetry-features}.

\begin{table}[H]
\centering
\caption{Bộ 8 đặc trưng telemetry chính thu thập từ Edge Node}
\label{tab:telemetry-features}
\begin{tabular}{|c|l|c|l|}
\hline
\textbf{STT} & \textbf{Feature} & \textbf{Đơn vị} & \textbf{Đại diện cho} \\ \hline
1 & \texttt{temperature\_c} & $^\circ\text{C}$ & Trạng thái nhiệt độ / Phần cứng (\textit{Thermal/Hardware}) \\ \hline
2 & \texttt{inference\_fps} & FPS & Thông lượng xử lý (\textit{Throughput}) \\ \hline
3 & \texttt{inference\_latency\_ms} & ms & Độ trễ suy luận (\textit{Inference performance}) \\ \hline
4 & \texttt{gpu\_usage\_pct} & \% & Áp lực tính toán GPU (\textit{GPU pressure}) \\ \hline
5 & \texttt{cpu\_usage\_pct} & \% & Tải tính toán CPU (\textit{CPU pressure}) \\ \hline
6 & \texttt{memory\_usage\_pct} & \% & Áp lực bộ nhớ RAM (\textit{Memory pressure}) \\ \hline
7 & \texttt{queue\_size} & gói/khung & Ứ đọng hàng đợi pipeline (\textit{Pipeline backlog}) \\ \hline
8 & \texttt{confidence\_mean} & $0 - 1$ & Chất lượng đầu ra mô hình (\textit{Model output quality}) \\ \hline
\end{tabular}
\end{table}

\noindent\textbf{Ý nghĩa đặc biệt của \texttt{confidence\_mean}:} 

Trong khi 7 đặc trưng đầu phản ánh hiệu năng và tải hệ thống/runtime, \texttt{confidence\_mean} là telemetry đại diện trực tiếp cho chất lượng suy luận của mô hình AI. Việc tích hợp đặc trưng này mang lại ý nghĩa quan trọng cho bài toán Edge AI:
\begin{itemize}
  \item Hệ thống có khả năng phát hiện hiện tượng suy giảm chất lượng dữ liệu đầu vào hoặc thoái hóa mô hình (\textit{input/model degradation}) ngay cả khi các chỉ số phần cứng (CPU, GPU, nhiệt độ) và hiệu năng (FPS, Latency) vẫn hoạt động ở mức hoàn toàn bình thường.
  \item Giúp phát hiện sớm hiện tượng camera bị mờ, góc quay lệch, ánh sáng thay đổi đột ngột hoặc camera bị che khuất mà các giải pháp giám sát hạ tầng truyền thống hoàn toàn bỏ sót.
\end{itemize}

\subsection{Kiến trúc phát hiện sự cố đa tầng và Feature Engineering}

Luồng thu thập, xử lý chuỗi thời gian và phát hiện sự cố được minh họa trong Hình~\ref{fig:monitoring-incident-flow}.

\begin{figure}[H]
\centering
\begin{tikzpicture}[node distance=0.55cm,
  nodebox/.style={draw, rounded corners, align=center, minimum height=0.75cm, text width=2.6cm, font=\small, fill=blue!4},
  subbox/.style={draw, rounded corners, align=center, minimum height=0.65cm, font=\footnotesize, fill=white},
  outerbox/.style={draw, dashed, rounded corners, inner sep=6pt, fill=gray!4}
]
  % Cột 1: Thu thập Telemetry
  \node[nodebox, text width=12cm] (edge) {\textbf{EDGE NODE TELEMETRY SOURCE}};
  
  \node[subbox, below left=0.4cm and 2.2cm of edge, text width=3.3cm] (hw) {\textbf{Hardware}\\Temp, CPU, GPU};
  \node[subbox, below=0.4cm of edge, text width=3.3cm] (rt) {\textbf{Runtime}\\FPS, Latency, Queue};
  \node[subbox, below right=0.4cm and 2.2cm of edge, text width=3.3cm] (ai) {\textbf{AI Output}\\Confidence Mean};

  \draw[flowarrow] (edge.south -| hw.north) -- (hw.north);
  \draw[flowarrow] (edge.south) -- (rt.north);
  \draw[flowarrow] (edge.south -| ai.north) -- (ai.north);

  % Cột 2: Sliding window
  \node[nodebox, below=1.8cm of edge, text width=12cm, fill=orange!6] (window) {\textbf{Sliding Window Feature Engineering (Cửa sổ 5 phút)}\\Trích xuất: \texttt{current}, \texttt{mean}, \texttt{min/max}, \texttt{std}, \texttt{slope}, \texttt{delta}};

  \draw[flowarrow] (hw.south) -- (window.north -| hw.south);
  \draw[flowarrow] (rt.south) -- (window.north);
  \draw[flowarrow] (ai.south) -- (window.north -| ai.south);

  % Cột 3: ML Classifier
  \node[nodebox, below=0.5cm of window, text width=8cm, fill=green!6] (ml) {\textbf{ML Classifier (XGBoost)}\\Phát hiện suy giảm đồng thời (\textit{Pipeline Degradation})};
  \draw[flowarrow] (window) -- (ml);

  % Cột 4: Output states
  \node[subbox, below left=0.5cm and 0.8cm of ml, text width=2.8cm, fill=green!10] (normal) {\textbf{NORMAL}\\Hệ thống ổn định};
  \node[subbox, below right=0.5cm and 0.8cm of ml, text width=2.8cm, fill=red!10] (warning) {\textbf{WARNING}\\Kích hoạt Incident};
  \node[nodebox, below=0.5cm of warning, text width=3.2cm, fill=purple!8] (llm) {\textbf{LLM AI Agent}\\Chẩn đoán \& Runbook};

  \draw[flowarrow] (ml.south) -| (normal.north);
  \draw[flowarrow] (ml.south) -| (warning.north);
  \draw[flowarrow] (warning) -- (llm);

\end{tikzpicture}
\caption{Kiến trúc xử lý đặc trưng Telemetry và phân loại sự cố kết hợp LLM Agent}
\label{fig:monitoring-incident-flow}
\end{figure}

Để mô hình học máy phân loại chính xác mà không bị nhiễu bởi các biến động tức thời, hệ thống không huấn luyện trực tiếp trên 8 giá trị đo tức thời đơn lẻ. Thay vào đó, áp dụng kỹ thuật \textbf{Sliding Window 5 phút} để trích xuất các đặc trưng thời gian (\textit{Temporal Features}):
\begin{itemize}
  \item \textbf{Giá trị tức thời \& trung bình:} \texttt{current}, \texttt{mean}.
  \item \textbf{Khoảng biến thiên \& độ phân tán:} \texttt{min}, \texttt{max}, \texttt{std} (độ lệch chuẩn).
  \item \textbf{Xu hướng biến thiên:} \texttt{slope} (độ dốc hồi quy), \texttt{delta} (mức thay đổi so với đầu cửa sổ).
\end{itemize}

\subsection{Cơ chế phát hiện suy giảm đồng thời (Pipeline Degradation)}

Điểm ưu việt của mô hình XGBoost phân loại trạng thái nằm ở khả năng học quan hệ tương quan phi tuyến giữa các chiều đặc trưng, phát hiện sự suy thoái toàn diện của pipeline Edge AI:
\begin{itemize}
  \item \textbf{Trường hợp suy thoái hệ thống (WARNING):} Khi $\text{GPU} \uparrow$, $\text{Temperature} \uparrow$, $\text{Latency} \uparrow$, $\text{Queue} \uparrow$ trong khi $\text{FPS} \downarrow$ và $\text{Confidence} \downarrow \downarrow$. Dù từng chỉ số có thể chưa chạm ngưỡng ngắt cứng (hard threshold), mô hình vẫn kịp thời dự đoán cảnh báo nguy cơ tắc nghẽn hoặc nghẽn cổ chai cục bộ.
  \item \textbf{Trường hợp tải tăng bình thường (NORMAL):} Khi $\text{GPU} \uparrow$ do có đối tượng xuất hiện liên tục trong khung hình, nhưng $\text{Temperature}$, $\text{Latency}$, $\text{FPS}$, $\text{Queue}$ và $\text{Confidence}$ vẫn duy trì ổn định $\rightarrow$ Mô hình xác định là trạng thái hoạt động bình thường, loại bỏ triệt để cảnh báo sai (\textit{false positives}).
\end{itemize}

Khi mô hình dự báo trạng thái \textbf{WARNING}, hệ thống sẽ tạo một bản ghi \textit{Incident}, tổng hợp snapshot metric, lịch sử lệnh và chuyển tiếp cho \textbf{LLM AI Agent} để thực hiện chẩn đoán nguyên nhân gốc và đề xuất giải pháp xử lý (\textit{Runbook}).